\documentclass[11pt,a4paper]{article}
\usepackage[utf8]{inputenc}
\usepackage[T1]{fontenc}
\usepackage[portuguese]{babel}
\usepackage{amsmath, amssymb, graphicx, float, setspace, hyperref}
\usepackage[top=1.5cm,bottom=1.5cm,left=1.75cm,right=1.75cm]{geometry}
\usepackage[tiny]{titlesec} 

\hypersetup{colorlinks=true, allcolors=blue}

\begin{document}
\onehalfspacing
\begin{center}
    \vspace*{0.5cm}
    {\LARGE \textbf{Trabalho 01: Room Reservation System}} \\
    \textit{Estrutura de Dados} \\
    \vspace{0.8cm}

    \begin{minipage}{0.45\textwidth}
        \centering
        \textbf{Rafael Gontijo Ferreira} \\
        \href{mailto:rafael.ferreira.2@fgv.edu.br}{\footnotesize \texttt{rafael.ferreira.2@fgv.edu.br}}
    \end{minipage}
    \begin{minipage}{0.45\textwidth}
        \centering
        \textbf{Rudá Dantas Ruoso Brandão} \\
        \href{mailto:rudabrandao@ruoso.com}{\footnotesize \texttt{rudabrandao@ruoso.com}}
    \end{minipage}
    
    \vspace{0.8cm}
\end{center}

\section*{Descrição do Projeto}
Este projeto apresenta a implementação de um sistema de reserva de salas, desenvolvido como parte da primeira atividade de Estrutura de Dados. O objetivo central é gerenciar a ocupação acadêmica de forma eficiente, respeitando restrições de horários e capacidade física.

O sistema permite as seguintes operações fundamentais:
\begin{itemize}
    \item \textbf{Reserva de salas:} Alocação de disciplinas com base em dia, horário e demanda de alunos.
    \item \textbf{Cancelamento:} Remoção de reservas existentes no sistema.
    \item \textbf{Visualização:} Exibição da grade de horários completa de cada sala.
\end{itemize}
Instrução de compilação:
\begin{verbatim}
    g++ main.cpp ReservationSystem.cpp ReservationRequest.cpp 
    -o Reservation_main_test.out 
    -std=c++23 -Wall -Wextra -Wconversion -Wpedantic -O2
\end{verbatim}
Instrução de execução:
\begin{verbatim}
    ./Reservation_main_test.out
\end{verbatim}
\section*{Organização do Sistema}
O projeto foi modularizado em classes para garantir a separação de responsabilidades e facilitar a manutenção:
\begin{itemize}
    \item \texttt{main.cpp}: Ponto de entrada que executa os testes automatizados e demonstra as funcionalidades.
    \item \texttt{ReservationRequest}: Define a estrutura da solicitação de reserva.
    \item \texttt{ReservationSystem}: Núcleo do sistema, responsável pelo gerenciamento dos dados e lógica.
\end{itemize}

\section*{Estruturas de Dados: Escolhas, Alternativas e Trade-offs}
Para cumprir os requisitos de não utilizar a \textit{Standard Template Library} (STL), o sistema foi construído sobre matrizes e alocação dinâmica de memória:

\begin{enumerate}
    \item \textbf{Matriz de Disponibilidade (Implementada):} Cada sala armazena uma matriz bidimensional estática de strings para $5$ dias e $14$ horários. \textbf{Por quê?} Essa estrutura permite acesso e verificação de um slot em tempo $O(1)$. 
    \item \textbf{Array Dinâmico de Salas (Implementado):} As salas são instanciadas dinamicamente e ordenadas no construtor por capacidade. Isso permite uma política de "Melhor Ajuste" ao iterar sequencialmente, alocando a menor sala viável para a disciplina.
    \item \textbf{Lista Encadeada (Alternativa Considerada):} Foi cogitado o uso de listas encadeadas para armazenar apenas as reservas ativas, substituindo a matriz densa por uma estrutura esparsa.
\end{enumerate}

\textbf{Trade-offs:} A matriz estática escolhida favorece a \textbf{velocidade e simplicidade} de implementação (acesso $O(1)$ a qualquer horário), mas possui o revés do \textbf{desperdício de memória}, já que todos os slots (mesmo vazios) são previamente alocados. Por outro lado, a alternativa da lista encadeada economizaria memória em ambientes com poucas reservas, mas sofreria um \textit{trade-off} de performance, degradando a busca para $O(K)$ (onde $K$ é o número de reservas na sala), além de aumentar o risco de vazamento de memória devido à manipulação complexa de ponteiros.

\section*{Custo Computacional}
Considerando $S$ como o número de salas, $D = 5$ (dias da semana) e $H = 14$ (horários diários):

\begin{itemize}
    \item \textbf{Verificação de conflito:} $O(H_{reserva})$. Está embutida na função de reserva. Para uma sala específica, checar se os blocos estão vazios exige iterar apenas pela duração solicitada na reserva (máximo de 14 iterações).
    \item \textbf{Reserva:} $O(S \times H_{reserva})$. O sistema varre o array de salas (previamente ordenado). No pior cenário, itera sobre todas as $S$ salas, avaliando a disponibilidade dos horários $H$ solicitados até encontrar uma vaga ou falhar, a complexidade assintótica reduz-se a $O(S)$.
    \item \textbf{Cancelamento:} $O(S \times D \times H)$. Como a solicitação de cancelamento recebe apenas o nome da disciplina, o algoritmo varre toda a matriz de ocupação de todas as salas para zerar as ocorrências. Sendo $D$ e $H$ constantes (70 slots), a complexidade assintótica reduz-se a $O(S)$.
    \item \textbf{Exibição da grade:} $O(S \times D \times H)$. Para imprimir a grade e agrupar horários contíguos (ex: 7h \textasciitilde 10h), o sistema iterativamente lê e compara todos os slots de memória do sistema. Da mesma forma, reduz-se assintoticamente a $O(S)$.
\end{itemize}

\end{document}